% Retained preamble for SIGMA 2023 article 008 cards, copied from
% sources/sigma-2023-008-an-askey-wilson-algebra-of-rank-2/pinned-original-source.tex
% (doi:10.3842/SIGMA.2023.008, SIGMA 19 (2023), CC BY 4.0).
%
% Omitted and why: the publisher \documentclass{sigma} block (the sigma
% class is unavailable in the toolchain; the assembler supplies the amsart
% class) and the publisher frontmatter commands between \begin{document} and
% the first section; no page geometry or font setup is carried so it cannot
% fight the forced xelatex engine. amsmath/amssymb/amsthm are supplied up front
% (the sigma class loaded them; amsthm is needed for \theoremstyle below) and
% graphicx arrives with the assembler line. Everything from the author's
% package block through the theorem declarations is copied verbatim from the
% pinned source (source lines 9-151),
% original line order and comment shape preserved: active packages kept live
% (phaistos is installed in the toolchain TeX tree; soul/xcolor/graphicx are
% reachable through the active lines), publisher-only or inactive lines keep
% their original commented shape so line numbers are preserved.
\usepackage{amsmath,amssymb,amsthm}
\usepackage{mathrsfs}
\usepackage{bm}

\numberwithin{equation}{section}

\newtheorem{Theorem}{Theorem}[section]
\newtheorem*{Theorem*}{Theorem}
\newtheorem{Corollary}[Theorem]{Corollary}
\newtheorem{Lemma}[Theorem]{Lemma}
\newtheorem{Proposition}[Theorem]{Proposition}
 { \theoremstyle{definition}
\newtheorem{Definition}[Theorem]{Definition}
\newtheorem{Note}[Theorem]{Note}
\newtheorem{Example}[Theorem]{Example}
\newtheorem{Remark}[Theorem]{Remark} }

\newcommand{\Dom}{\rm Dom}

\newcommand{\supp}{\operatorname{supp}}

\newcommand{\leaderfill}{\leaders\hbox to 1em{\hss.\hss}\hfill}
\newcommand{\RR}{\mathbb{R}}
\newcommand{\R}{\mathbb{R}}
\newcommand{\CC}{\mathbb{C}}
\newcommand{\QQ}{\mathbb{Q}}
\newcommand{\ZZ}{\mathbb{Z}}
\newcommand{\NN}{\mathbb{N}}
\newcommand{\LL}{\mathbb{L}}
\newcommand{\PP}{\mathbb{P}}
\renewcommand{\P}{{\mathbb P}}
\newcommand{\EE}{\mathbb{E}}
\newcommand{\E}{\mathbb{E}}
\newcommand{\FF}{\mathbb{F}}
\newcommand{\DD}{\mathbb{D}}
\newcommand{\D}{\mathbb{D}}
\newcommand{\e}{\varepsilon}
\newcommand{\one}{\mathbf{1}}

\newcommand{\calL}{{\mathscr L}}
\newcommand{\OO}{\Omega}
\newcommand{\A}{{\mathscr A}}
\newcommand{\F}{{\mathscr F}}
\newcommand{\G}{{\mathscr G}}
\renewcommand{\H}{{\mathscr H}}
\newcommand{\g}{\gamma}
\newcommand{\lb}{\langle}
\newcommand{\rb}{\rangle}
\newcommand{\limn}{\lim_{n\to\infty}}
\newcommand{\limk}{\lim_{k\to\infty}}
\newcommand{\umd}{\textsc{umd}}
\newcommand{\nn}{|\!|\!|}
\newcommand{\CIS}{CIS}
\newcommand{\AT}{AT}

\newcommand{\abcd}{\alpha,\beta,\gamma,\delta}
\newcommand{\aq}{\mathcal{A}_q(\mathrm{SU}(2))}
\newcommand{\aql}{\mathcal{A}_q(\mathrm{SL}(2,\CC))}
\newcommand{\aqll}[1]{\text{A}_q^{#1}}
\newcommand{\aw}{\mathrm{AW}\!\left(3\right)}
\newcommand{\awc}[2]{\mathrm{AW}_{#1}\!\left(#2\right)}
\newcommand{\AW}[1]{\mathrm{AW}_{#1}}
\newcommand{\awr}{\mathscr{A}\!\!\mathscr{W}}
\newcommand{\bivariateawalgebra}{rank 2 Askey--Wilson algebra }
\newcommand{\bigo}[1]{\mathcal{O}\left(#1\right)}
\newcommand{\chq}{\mathrm{cosh}_q}
\newcommand{\dif}[2]{\frac{\mathrm{d}#1}{\mathrm{d}#2}}
\newcommand{\dx}{\!\mathrm{d}x}
\newcommand{\dmu}[1][x]{\!\mathrm{d}\mu(#1)}
\newcommand{\dyone}{\Delta\left(Y_{L}\right)}
\newcommand{\dyzero}{\Delta\left(Y_{K}\right)}
\newcommand{\gen}{\mathscr{A}}
\newcommand{\germ}[1]{\mathfrak{#1}}
\newcommand{\id}{\mathrm{id}}
\newcommand{\inprod}[2]{\left\langle #1, #2 \right\rangle}
\newcommand{\half}{\frac{1}{2}}
\newcommand{\hyp}[5]{\,_{#1}F_{#2} \left[ \begin{matrix}
		#3 \\ #4 \end{matrix};#5 \right]}
\newcommand{\lp}{\lambda_p}
\newcommand{\ls}{\lambda_s}
\newcommand{\matrixa}{\left(\begin{matrix}
		a & b \\ c & d
	\end{matrix}\right)}
\newcommand{\matrixalpha}{\left(\begin{matrix}
		\alpha & \beta \\ \gamma & \delta
	\end{matrix}\right)}
\newcommand{\matrixg}{\left(\begin{matrix}
		g_{11} & g_{12} \\ g_{21} & g_{22}
	\end{matrix}\right)}
\newcommand{\matrixid}{\left(\begin{matrix}
		1 & 0 \\ 0 & 1
	\end{matrix}\right)}
\newcommand{\matrixzelf}[4]{\left(\begin{matrix}
		#1 & #2 \\ #3 & #4
	\end{matrix}\right)}
\newcommand{\multivariateawalgebra}{multivariate Askey--Wilson algebra }
\newcommand{\Multivariateawalgebra}{Multivariate Askey--Wilson algebra }
\newcommand{\ncr}[2]{\left[ \begin{matrix}
		#1 \\ #2 \end{matrix} \right]}
\newcommand{\qncr}[3][q]{\left[\begin{matrix}
		#2 \\ #3 \end{matrix} \right]_{#1}}
\newcommand{\para}{\mathscr{W}}
\newcommand{\psimn}[2]{\psi_{#1}^{#2}\otimes \psi_{#2}}
\newcommand{\polaq}{\text{Pol}\left(\aq\right)}
\newcommand{\polaql}{\text{Pol}\left(\aql\right)}
\newcommand{\polg}{\text{Pol}\left(G\right)}
\newcommand{\polsl}{\text{Pol}(\text{SL}(2,\CC))}
\newcommand{\qhyp}[5]{\,_{#1}\phi_{#2} \left[ \begin{matrix}
		#3 \\ #4 \end{matrix};#5 \right]}
\newcommand{\rhots}{\rho_{\tau,\sigma}}
\newcommand{\shq}{\mathrm{sinh}_q}
\newcommand{\slc}{\mathfrak{sl}(2,\CC)}
\newcommand{\SL}{\text{SL}(2,\CC)}
\newcommand{\SU}{\text{SU}(2)}
\makeatletter
\newcommand{\subalign}[1]{%
	\vcenter{%
		\Let@ \restore@math@cr \default@tag
		\baselineskip\fontdimen10 \scriptfont\tw@
		\advance\baselineskip\fontdimen12 \scriptfont\tw@
		\lineskip\thr@@\fontdimen8 \scriptfont\thr@@
		\lineskiplimit\lineskip
		\ialign{\hfil$\m@th\scriptstyle##$&$\m@th\scriptstyle{}##$\hfil\crcr
			#1\crcr
		}%
	}%
}
\makeatother
\newcommand{\sweed}[2]{#1_{(#2)}}
\newcommand{\tensor}[1]{^{\otimes{#1}}}
\newcommand{\tr}[1]{\text{tr}\left(#1\right)}
\newcommand{\ts}{\tau,\sigma}
\newcommand{\uq}{\mathcal{U}_q(\mathfrak{su}(2))}
\newcommand{\Uq}{\mathcal{U}_q}
\newcommand{\UqO}[1][q]{\mathcal{U}_{#1}\left(\Omega_0\right)}
\newcommand{\uqeen}{\mathcal{U}_q(\mathfrak{su}(1,\!1))}
\newcommand{\uql}[1][q]{\mathcal{U}_{#1}(\mathfrak{sl}(2,\CC))}
\newcommand{\uqlO}[1][q]{\mathcal{U}_{#1}\left(\mathfrak{sl}(2,\CC),\Omega_0\right)}
\newcommand{\usl}{\mathcal{U}(\mathfrak{sl}(2,\CC))}
\newcommand{\usu}{\mathcal{U}(\mathfrak{su}(2,\CC))}
\newcommand{\vectornorm}[1]{\left|\left|\left|#1\right|\right|\right|}
\newcommand{\vspan}[1]{\text{span}\{#1\}}
\newcommand{\yone}{Y_{L}\otimes 1}
\newcommand{\yzero}{1\otimes Y_{K}}

% --- lane-s11 citation neutralization stub (no-bibliography card convention; shim family x2 PARK-CITE-STUB):
% the author's proofs cite the source bibliography by key; this card ships no bibliography, so render the
% key (and optional locator) as a textual pointer instead of resolving it.
\renewcommand{\cite}[2][]{\mbox{[#2\if!#1!\else, #1\fi]}}
